% =============================================================================
% CHAPTER 1: Introduction
% =============================================================================
% SPDX-FileCopyrightText: 2026 Drewan Rutherford
% SPDX-License-Identifier: LPPL-1.3c
%
% Suggested sections: Motivation and context; Aims and objectives; Contributions;
% Dissertation structure (roadmap). Optional: Scope (boundaries of the work),
% Terminology or key terms. Start with a short chapter intro (no \section) before
% the first section. This file also demonstrates template features; see
% Section~\ref{sec:ch1-reference}. Remove that section for your own report.
% =============================================================================

% -----------------------------------------------------------------------------
% Chapter intro (no section): one short paragraph before the first \section
% -----------------------------------------------------------------------------
This chapter introduces the dissertation. It outlines the motivation and context for the work, the aims and objectives, the main contributions, and the structure of the remaining chapters. The following sections expand on each of these.

The rest this chapter details the \nameref{sec:introduction/motivation_and_context}, as well as the \nameref{sec:introduction/aims_and_objectives} of the project.

Following this chapter~\ref{ch2:background}, will further flesh out the context, introduce and compare existing games and programs alongside their strengths and weaknesses.

Chapter~\ref{ch3:method} explains the methodology of both the development and evaluation of the Code Heist prototype.

Chapter~\ref{ch4:requirements} will detail the requirements used to develop the prototype and their reasoning.

Chapter~\ref{ch5:design} describes and justifies the design of the prototype. The prototype produced for this dissertation was far more complex than anticipated, but -- potentially, as a result of this -- serves as an excellent blueprint or foundation for future projects.

Chapter~\ref{ch6:results} critically analyses the prototype using the methods described in chapter 3, as well as providing a subjective evaluation of the prototype from the developer as well as approach itself.

Finally, chapter~\ref{ch7:conclusions} details conclusions that can be drawn from the successes and shortcomings of this project. It will also go into detail about future work that can be spun off from this dissertation including alternative game genres, future methods of evaluation and potential future of Code Heist. There is also some discussion about the cut content from the prototype, although some of this had to be pushed into [TODO: list appendix].

Beyond this, chapter~\ref{ch8:reflections} includes some reflections on the project as a whole from the author.

As mentioned in frontmatter under \nameref{sec:source_code}, the source code is available on GitHub in both its modern state as well as its state at time of this report's submission.

\section{Motivation and Context}\label{sec:introduction/motivation_and_context}

\section{Aims and Objectives}\label{sec:introduction/aims_and_objectives}

The aim of this project was to evaluate the efficacy of using video-games to teach programming through its games mechanics whilst avoiding the aforementioned ‘instruction sequence’ paradigm. This was be done by developing simulating a short development period to make a platforming game and then evaluating the results.

\textbf{The project had the following objectives:}

\begin{enumerate}[label=\textbf{Objective \arabic*:}, leftmargin=*]
    \item Research the Welsh curriculum (and English if required as a backup) and design the progression of the game around it, incorporating required programming techniques such as iteration, into its mechanics.
    \item Attempt to design and program a game that incorporates programming into its mechanics and core game-loop, that avoids the ‘instruction sequence’ paradigm.
    \item Evaluate the plausibility of designing said game within the tight time-frame presented by this project.
    \item If possible, evaluate the game to measure how approachable and enjoyable users find the experience using user testing and surveys, accounting for experience with both video games and programming.
\end{enumerate}\label{objectives}

\section{Contributions}\label{sec:introduction/contributions}

All of the code produced for this project was written by hand or through GUI tools within the Godot game engine's software for said engine. This project would not have been possible without this engine.